% !TeX root = ../main.tex
\section{Conclusion}
\label{sec:conclusion}

We presented a behavioral simulation framework for analog compute-in-memory inference that disentangles algorithmic generalization failures from physical precision-retention constraints. By comparing ideal pure-quantization baselines against realistic phase-change memory (PCM) physics, the framework serves as a practical decision aid for establishing deployment frontiers in low-precision analog hardware.

Several observations stand out. Under the tested PCM UnitCell simulation regime, convergence is observed across the 4/6/8-bit precision ladder where pure quantization fails. This convergence exposes a non-monotonic precision--accuracy curve: 8-bit PCM achieves the highest fresh accuracy (77.60\%) with negligible drift, 4-bit PCM achieves comparable fresh accuracy (76.68\%) at the cost of a 4.01~pp one-day drift penalty, and 6-bit PCM enters a D2D-sensitive transition zone with lower mean accuracy (68.55\%, seed-level std 6.03\%) but drift stability. The 6-bit dip reflects a quantization-D2D interaction where, at 1/64 resolution, device mismatch perturbations become comparable to quantization steps. Separately, in the ideal-device ablation, fixed hardware-instance transfer is the most prominent algorithmic risk observed in our matrix: naive 4-bit deployment collapses to 14.64\%, whereas Ensemble Hardware-Aware Training (HAT) resamples mismatch masks to recover fresh-instance accuracy to 86.16$\pm$0.19\%.

Ultimately, this framework clarifies that standard fixed-mask training is insufficient for extreme low-precision analog deployment. Ensemble HAT resolves this algorithmic bottleneck, allowing system designers to navigate the subsequent physical limits of the analog precision-retention frontier. Our findings suggest that noise-aware training principles are foundational for memory-bound inference components, including analog KV-cache architectures.
